\section{Materials and Methods}
\subsection{Materials}
Analysis of the interplay between \acp{drm}, and IEMs were performed upon datasets upon the \ac{lahivdb} and the Stanford University HIV Drug Resistance Database \ac{hivdrmdb}. NetMHCPan version 4.1 and an in-house python script were used in combination for accommodating large data throughput \cite{reynissonNetMHCpan4NetMHCIIpan4Improved2020}. Within the Python script, SciPy and Pandas data processing libraries were used for the analysis \cite{2020SciPy-NMeth,the_pandas_development_team_2024_10957263}. Figures were generated with use of Matplotlib and Seaborn \cite{Hunter:2007,Waskom2021}.

The \ac{lahivdb} dataset of known HIV-1 \ac{ctl} epitopes including the peptide sequence and position relative to the gene. Datasets were individually downloaded by querying for integrase, p24, protease, and \ac{rt} while all other search parameters were left as the most broad casting option i.e., selecting the "ALL" option when present. The \ac{hivdrmdb} amongst other tools, and other datasets, contains a scoring mechanism for known drug resistance inducing mutations in the HIV-1 genome \cite{rheeHIV1TransmittedDrug2020}. Spreadsheets of such scoring for the different drugs with respect to a given single mutation were downloaded. NetMHCPan version 4.1 was used to predict binding against different \ac{hla} alleles \cite{reynissonNetMHCpan4NetMHCIIpan4Improved2020}. Data ingress and egress, manipulation and vectorized operations were performed by Pandas and the SciPy library was used for statistical associations \cite{the_pandas_development_team_2024_10957263,2020SciPy-NMeth}. Lastly, Seaborn was used for the automated generation and regression figures by making use of Matplotlib \cite{Waskom2021,Hunter:2007}. 

\subsection{Methods}

All aforementioned datasets were imported from respective \ac{csv} sheets downloaded from their respective databases. This was then followed by the generation of new epitopes based off the set of known epitopes, and the set of drug resistant mutations. All epitopes, natural or generated were then run through NetMHCPan 4.1 to predict their binding level towards various \ac{hla} alleles. Such binding levels were fit into categories. Hypothetical epitopes were generated from the \acp{naepi}. All generated epitopes were cleaned and filtered to increase the number of interpretable natural epitopes and reduce reduce computation time respectively before running the dataset in NetMHCPan \cite{reynissonNetMHCpan4NetMHCIIpan4Improved2020}. \Acp{drm} and their associated resistance scores were summed in accordance to the guidelines of the \ac{hivdrmdb}. The resistance scores and \ac{drm} scores were then categorized according to their respective guidelines from \ac{hivdrmdb}. Association between the presence of \acp{drm} and the binding level with various \ac{hla} alleles were tested and p-values were aggregated.

The first step involved downloading the data and parsing the data into Python. Data was procured from both \ac{lahivdb} and \ac{hivdrmdb} as described in the materials section. \Ac{drm} were further parsed and merged into a large, master spreadsheet where the drugs irrelevant to a mutation were set to have a \ac{drm} of 0. Similarly, known \acp{naepi} were also parsed. Unused columns of metadata were dropped and epitopes with the exact identical sequences were dropped immediately as to ensure that any associations found would be purely based on their occurrence in a epitope, and not biased by the frequency in a population. Furthermore, the \ac{lahivdb} contained synthetic epitopes created by concatenating two epitopes together. Such epitopes did not have well-defined bounds and thus, also did not have well-defined \acp{drm} \cite{mamroshHIVMolecularImmunology2022}. Such junction-based epitopes had their junctions removed and the original epitope restored by selecting the portion of the junction epitope based on the annotated bound.

The unique set of naturally occurring HIV epitopes then had variants generated naively simply by generating all possible different combinations of mutations that could be added to a naturally occurring epitope was created. These epitopes were then added to a new flat database with their original \acp{drm}, i.e., mutations that existed before any modification, and the added mutations. Notably, due to some epitopes having many possible combination of \acp{drm}, constraints were added such that generated epitopes that resulted in a \ac{drm} score lower than 15 was removed. This corresponds to \acp{hivdrmdb} categorization of having none, or potentially a very low level drug resistance (Table \ref{table:hivdrmdb_categorizations}). Additionally, some \ac{naepi} had many peptide positions with known \acp{drm} and thus, a ranked threshold was applied where the \ac{drm} scores were used to rank the generated epitopes and such epitopes that ranked higher than 1000 were removed. I.e., the top 1000 highest \ac{drm} scoring generated epitopes were kept.

Then, the produced aggregate of both \acp{naepi} and the generated epitopes had binding levels with all \ac{hla} alleles that were recognized by NetMHCPan (table \ref{table:NetMHCPan_supported_hla}) predicted and cataloged \cite{reynissonNetMHCpan4NetMHCIIpan4Improved2020}. The resulting database had \ac{drm} scores and \ac{pelrank} calculated and mapped according to table \ref{table:hivdrmdb_categorizations} and to table \ref{table:elrankpercentile_categorizations} respectively (for more information on this metric, see original software paper for version 4.0) \cite{reynissonNetMHCpan4NetMHCIIpan4Improved2020,jurtzNetMHCpanImprovedPeptideMHC2017}. Furthermore, changes in binding level were calculated by comparing the \acp{naepi} with the generated ones. 

\begin{table}
    \centering % TODO decide if this is the right command.
    \begin{tabular}{||c||}
        \hline
        Supported HLA Variants \\
        \hline
        \hline
        HLA-A02:01 \\
        HLA-A01:01 \\
        HLA-A03:01 \\
        HLA-A24:02 \\
        HLA-A26:01 \\
        HLA-B07:02 \\
        HLA-B08:01 \\
        HLA-B27:05 \\
        HLA-B39:01 \\
        HLA-B40:01 \\
        HLA-B58:01 \\
        HLA-B15:01 \\
        \hline
    \end{tabular}
    \caption{The supported list of HLA alleles for which NetMHCPan is capable of predicting \ac{pelrank} \cite{reynissonNetMHCpan4NetMHCIIpan4Improved2020}.}
    \label{table:NetMHCPan_supported_hla}
\end{table}
\begin{table}
    \centering % TODO decide if this is the right command.
    \begin{tabular}{||c|c||}
        \hline
        Resistance Category & \Ac{drm} Score (Inclusive) \\
        \hline
        \hline
        No drug resistance & 0 to 9 \\
        \hline
        Potential low-level drug resistance & 10 to 14 \\
        \hline
        Low-level drug resistance & 15 to 29 \\
        \hline
        Intermediate resistance & 30 to 59 \\
        \hline
        High-level resistance & 60 and above \\
        \hline
    \end{tabular}
    \caption{The categorical levels of drug resistance based on the \ac{drm} score assigned for an entire sequence as defined by \ac{hivdrmdb}.}.
    \label{table:hivdrmdb_categorizations}
\end{table}
\begin{table}[h]
    \centering % TODO decide if this is the right command.
    \begin{tabular}{||c|c||}
        \hline
        Resistance Category & \Ac{drm} Score (Inclusive) \\
        \hline
        \hline
        Strong Binding & Less than 0.5 \\
        \hline
        Weak Binding & Less than 2, greater than or equal to 0.5 \\
        \hline
        No Binding & Other \\
        \hline
    \end{tabular}
    \caption{The categorical levels of HLA binding based on the percentile rank predicted by NetMHCPan 4.1 where both the metric and the thresholds were advised by the original paper and authors of the software \cite{reynissonNetMHCpan4NetMHCIIpan4Improved2020}.}
    \label{table:elrankpercentile_categorizations}
\end{table}

Lastly, all combinations between all sets of \acp{drm}, and the empty set of \acp{drm}, i.e., the lack of \acp{drm}, were tested for association with all different levels of binding, that is to say, these combinations of \acp{drm} were tested against the occurrences of weak, strong, and no binding level epitopes. The Fisher's exact test was used due to the relatively small sample count of epitopes with the same set of \acp{drm} with unique epitopes. This process generated a list of different p-values for each possible combination. This list of p-values was then filtered by a $\alpha=0.05$.